% Figure for Relativity §C10.4, Example §C10.4.2 and Theorem §C10.4.2: (a) Minkowski space with the origin removed,
% the worldline γ(t) = (t, 0), t < 0, has no endpoint; (b) the Rindler wedge x > |t| with ξ = const hyperbolae and η = const rays,
% the inertial worldline x = x0 = 1 inside it for proper time 2x0. Generated by rel-C/tools/figs/c10_4_figs.py.
\begin{tikzpicture}[font=\small, scale=1.1]
\draw[->, black!55] (-1.9,0) -- (1.9,0) node[right] {$x$};
\draw[->, black!55] (0,-1.9) -- (0,1.9) node[above] {$t$};
\draw[black!35, dashed] (-1.15,-1.15) -- (1.6,1.6); \draw[black!35, dashed] (-1.6,1.6) -- (1.15,-1.15);
\draw[figB, very thick, -{Stealth[length=5pt]}] (0,-1.7) -- (0,-0.55);
\draw[figB, very thick] (0,-0.6) -- (0,-0.07);
\draw[black!45, thick] (0.45,-1.7) -- (0.45,1.7);
\node[black!60, font=\scriptsize, anchor=west] at (0.48,1.45) {passes by: complete};
\draw[black, fill=white] (0,0) circle (2.2pt);
\node[font=\scriptsize, anchor=south] at (0,0.1) {removed};
\node[figB, font=\scriptsize, anchor=east, align=right] at (-0.2,-1.5) {$\gamma(t) = (t, 0)$\\$t < 0$: ends\\after $\tau = |t_0|$};
\node at (0,-2.35) {(a) $\mathbf R^{1,1}\setminus\{0\}$};
\begin{scope}[shift={(4.6,0)}]
\fill[figO!10] (0,0) -- (1.9,1.9) -- (1.9,-1.9) -- cycle;
\draw[->, black!55] (-0.6,0) -- (2.1,0) node[right] {$x$};
\draw[->, black!55] (0,-1.9) -- (0,1.9) node[above] {$t$};
\draw[figO, thick, dashed] (0,0) -- (1.9,1.9) node[above, font=\scriptsize] {$t = x$};
\draw[figO, thick, dashed] (0,0) -- (1.9,-1.9);
\draw[figO!70] (1.289,-1.188) -- (1.242,-1.137) -- (1.198,-1.088) -- (1.155,-1.041) -- (1.114,-0.996) -- (1.075,-0.952) -- (1.038,-0.910) -- (1.003,-0.869) -- (0.969,-0.830) -- (0.936,-0.792) -- (0.905,-0.755) -- (0.876,-0.719) -- (0.848,-0.685) -- (0.821,-0.651) -- (0.796,-0.619) -- (0.772,-0.588) -- (0.749,-0.557) -- (0.727,-0.528) -- (0.706,-0.499) -- (0.687,-0.471) -- (0.669,-0.444) -- (0.651,-0.418) -- (0.635,-0.392) -- (0.620,-0.367) -- (0.606,-0.342) -- (0.593,-0.318) -- (0.580,-0.295) -- (0.569,-0.272) -- (0.559,-0.249) -- (0.549,-0.227) -- (0.541,-0.205) -- (0.533,-0.184) -- (0.526,-0.163) -- (0.520,-0.142) -- (0.514,-0.121) -- (0.510,-0.101) -- (0.506,-0.080) -- (0.504,-0.060) -- (0.502,-0.040) -- (0.500,-0.020) -- (0.500,0.000) -- (0.500,0.020) -- (0.502,0.040) -- (0.504,0.060) -- (0.506,0.080) -- (0.510,0.101) -- (0.514,0.121) -- (0.520,0.142) -- (0.526,0.163) -- (0.533,0.184) -- (0.541,0.205) -- (0.549,0.227) -- (0.559,0.249) -- (0.569,0.272) -- (0.580,0.295) -- (0.593,0.318) -- (0.606,0.342) -- (0.620,0.367) -- (0.635,0.392) -- (0.651,0.418) -- (0.669,0.444) -- (0.687,0.471) -- (0.706,0.499) -- (0.727,0.528) -- (0.749,0.557) -- (0.772,0.588) -- (0.796,0.619) -- (0.821,0.651) -- (0.848,0.685) -- (0.876,0.719) -- (0.905,0.755) -- (0.936,0.792) -- (0.969,0.830) -- (1.003,0.869) -- (1.038,0.910) -- (1.075,0.952) -- (1.114,0.996) -- (1.155,1.041) -- (1.198,1.088) -- (1.242,1.137) -- (1.289,1.188);
\draw[figO!70] (1.906,-1.176) -- (1.860,-1.100) -- (1.818,-1.027) -- (1.778,-0.955) -- (1.741,-0.885) -- (1.707,-0.816) -- (1.676,-0.748) -- (1.648,-0.682) -- (1.622,-0.616) -- (1.598,-0.552) -- (1.577,-0.488) -- (1.559,-0.426) -- (1.543,-0.363) -- (1.530,-0.302) -- (1.519,-0.241) -- (1.511,-0.180) -- (1.505,-0.120) -- (1.501,-0.060) -- (1.500,0.000) -- (1.501,0.060) -- (1.505,0.120) -- (1.511,0.180) -- (1.519,0.241) -- (1.530,0.302) -- (1.543,0.363) -- (1.559,0.426) -- (1.577,0.488) -- (1.598,0.552) -- (1.622,0.616) -- (1.648,0.682) -- (1.676,0.748) -- (1.707,0.816) -- (1.741,0.885) -- (1.778,0.955) -- (1.818,1.027) -- (1.860,1.100) -- (1.906,1.176);
\draw[figO!45] (0,0) -- (1.900,-1.020);
\draw[figO!45] (0,0) -- (1.900,1.020);
\draw[black!45, densely dotted, thick] (1.0,-1.85) -- (1.0,-1.0);
\draw[black!45, densely dotted, thick] (1.0,1.0) -- (1.0,1.85);
\draw[figB, very thick, -{Stealth[length=5pt]}] (1.0,-1.0) -- (1.0,0.3);
\draw[figB, very thick] (1.0,0.25) -- (1.0,1.0);
\fill[figB] (1.0,1.0) circle (1.6pt); \fill[figB] (1.0,-1.0) circle (1.6pt);
\node[figB, font=\scriptsize, anchor=west, align=left] at (2.0,1.0) {$x = x_0$: inside\\for $\tau = 2x_0$};
\node[figO, font=\scriptsize, anchor=north west] at (1.5,-0.02) {$\xi = 1.5$};
\node[figO, font=\scriptsize, anchor=north west] at (0.5,-0.02) {$0.5$};
\node at (0.7,-2.35) {(b) Rindler wedge $x > |t|$};
\end{scope}
\end{tikzpicture}
